% F085｜真实数学图；代表函数不构成抽象解的唯一性证明。
\begin{center}
\begin{tikzpicture}
\begin{axis}[width=12.4cm,height=5.1cm,axis lines=middle,
 xmin=-0.3,xmax=8.9,ymin=-2.7,ymax=3.8,
 xtick={1,2,4,8},xticklabels=\empty,ytick={-2,-1,0,1,2,3},xlabel={$t$},ylabel={$f(t)$},
 title={对数曲线 $f(t)=\log_2t$},title style={font=\small,yshift=3pt},
 tick label style={font=\small},label style={font=\small},
 ylabel style={at={(axis description cs:0,1)},anchor=south east,xshift=-18pt},
 clip=false,grid=major,grid style={black!8},enlargelimits=false,
 every axis plot/.append style={thick},samples=161]
\addplot[CExplanation,very thick,domain=0.2:8.3] {ln(x)/ln(2)};
\draw[dashed,black!35] (axis cs:2,0) |- (axis cs:0,1);
\draw[dashed,black!35] (axis cs:4,0) |- (axis cs:0,2);
\draw[dashed,black!35] (axis cs:8,0) |- (axis cs:0,3);
\addplot[only marks,mark=*,mark size=2pt,orange!80!black] coordinates {(2,1) (4,2) (8,3)};

\node[font=\small,anchor=north,yshift=-3pt] at (axis cs:1,-2.7) {$1$};
\node[font=\small,anchor=north,yshift=-3pt] at (axis cs:2,-2.7) {$2$};
\node[font=\small,anchor=north,yshift=-3pt] at (axis cs:4,-2.7) {$4$};
\node[font=\small,anchor=north,yshift=-3pt] at (axis cs:8,-2.7) {$8$};

\end{axis}
\end{tikzpicture}
\par\small 输入从 $2$ 到 $4$ 再到 $8$，每次乘 $2$；高度从 $1$ 到 $2$ 再到 $3$，每次加 $1$。
\end{center}
